The reference room-temperature structure contains five Cu sites with multiplicities in the ratio $1:1:1:2:2$ \citep{Koto1970}. Expanding the published coordinates and symmetry operations gives the environments in Table~\ref{tab:structural}. Cu1 and Cu3 are tetrahedrally coordinated, whereas Cu2, Cu4, and Cu5 have three nearest sulfur neighbors. The fourth neighbor of Cu4 is substantially farther away, at approximately 2.94~\AA. The resulting count of eight tetrahedrally and twenty triangularly coordinated Cu atoms per cell agrees with the crystallochemical comparison of Evans \citep{Evans1981}.

Thus, the four high-frequency NQR groups cannot all be assigned to four distinct triangular sites of unchanged $Pnma$ anilite: that reference structure has only three such sites. If each group corresponds to a different site of this phase, at least one must involve a tetrahedral environment. A change of low-temperature symmetry or contributions from another phase would require a different counting argument.

Cu1--Cu3 occupy Wyckoff position $4c$ with mirror site symmetry $.m.$; Cu4 and Cu5 occupy general positions $8d$ with site symmetry $1$. None of these symmetries requires $\eta=0$. At a mirror site, invariance under reflection perpendicular to $b$ gives $V_{ab}=V_{bc}=0$. The $b$ direction is therefore an EFG principal direction, but symmetry neither selects it as the direction of the largest principal value nor requires the other two eigenvalues to be equal.

These multiplicities are population constraints, not automatically measured spin-echo intensity ratios. Relaxation, excitation bandwidth, pulse efficiency, overlap, and instrumental frequency response affect the observable amplitudes.

X-ray diffraction supplies atomic coordinates, occupancies, displacement parameters, and symmetry. Numerical EFG asymmetries additionally require the local charge distribution. They cannot generally be identified with a bond-length distortion or an atomic displacement parameter. The structural reference used here is the near-room-temperature coordinate set archived as AMCSD~0009376, rather than a verified diffraction refinement of the present material at 5~K.

\subsection{Exploratory electrostatic model}

To assess what the reference geometry might suggest, a separate point-charge calculation was performed without using the NMR fit parameters. Define the geometrical tensor
\begin{equation}
G_{\alpha\beta}
=\sum_{j,\bm R}'q_j
\frac{3r_{j,\alpha}r_{j,\beta}
-r_j^2\delta_{\alpha\beta}}{r_j^5},
\label{eq:pointcharge}
\end{equation}
where the sum includes periodic images, $q_j$ is expressed in units of the elementary charge, and the observation-site self term is omitted. A Coulomb prefactor would convert $G$ to an electrostatic field gradient, but the full nuclear EFG also contains electronic contributions not represented in this model.

The baseline model assigns $q_{\Cu}=8/7$ and $q_{\mathrm S}=-2$, ensuring neutrality of the 28-Cu, 16-S unit cell. These are trial charges, not an experimental oxidation-state assignment. Periodic sums were evaluated by the Ewald method and checked by varying the splitting parameter and comparing with direct finite-radius sums. Additional neutral scenarios start from $q_{\Cu}=1$ and assign the remaining four positive charge units to one Cu orbit at a time: the charge increment on that orbit is $4/m_s$, where $m_s$ is its unit-cell multiplicity. All coordinates are held fixed.

\begin{table*}[t]
\centering
\small
\setlength{\tabcolsep}{4pt}
\caption{Crystallographic reference and exploratory point-charge results. Site labels and coordinates follow the room-temperature Koto--Morimoto structure \citep{Koto1970}; coordination and the numerical columns were calculated here. $\eta_{\rm pc}$ describes $G$; it is not an experimental determination or an electronic-structure prediction of the nuclear-EFG asymmetry. The relative magnitude is normalized to Cu2 in the uniform-Cu-charge model. The final column spans six chosen neutral charge scenarios and is not an experimental uncertainty interval.}
\label{tab:structural}
\begin{tabular}{llrlrrr}
\toprule
Site & Wyckoff / symmetry & $m_s$ & Nearest-S coordination
& $\eta_{\rm pc}$ & $|G_{ZZ}|/|G_{ZZ}^{\Cu2}|$
& Scenario span of $\eta_{\rm pc}$ \\
\midrule
Cu1 & $4c$ / $.m.$ & 4 & Distorted tetrahedral & 0.557 & 0.688 & 0.095--0.774 \\
Cu2 & $4c$ / $.m.$ & 4 & Nearly regular triangular & 0.059 & 1.000 & 0.018--0.165 \\
Cu3 & $4c$ / $.m.$ & 4 & Tetrahedral & 0.255 & 0.088 & 0.069--0.842 \\
Cu4 & $8d$ / $1$ & 8 & Triangular & 0.163 & 0.847 & 0.035--0.537 \\
Cu5 & $8d$ / $1$ & 8 & Distorted triangular & 0.585 & 0.951 & 0.566--0.642 \\
\bottomrule
\end{tabular}
\end{table*}

The baseline values in Table~\ref{tab:structural} should be read as hypotheses generated by a specified electrostatic model. They omit the on-site electronic charge distribution, covalency, screening, and antishielding. Numerical convergence of the lattice sum does not remove this physical limitation, and its magnitude is not converted into predicted NQR frequencies.

Two tendencies are useful for designing an assignment test. First, Cu2 has a nearly regular triangular environment and a small model asymmetry. It is a plausible candidate for a nearly axial contribution, but matching this single number is insufficient: Cu2 also has one of the largest model $|G_{ZZ}|$, whereas the near-axial fitted group s2 is not the highest-frequency NQR group. Frequencies and asymmetries must be matched together.

Second, Cu3 has a much smaller model tensor magnitude than the other four sites. It remains the smallest in all six charge scenarios, motivating a possible low-EFG contribution from this tetrahedral site. This is consistent with investigating an overlapping or low-frequency fifth response, but does not identify an observed resonance or establish a frequency for it. A nearly regular tetrahedron should not be equated with $\eta\simeq0$: exact tetrahedral symmetry makes the entire EFG vanish, whereas a small residual tensor can have a sizeable and model-sensitive asymmetry.

The sensitivity is substantial even within this elementary approach. For Cu1, the nearest sulfur shell alone gives $\eta_{\rm pc}\simeq0.021$, while the complete periodic baseline model gives 0.557. Geometry of the first coordination shell is therefore insufficient even before electronic EFG contributions are considered. A quantitative crystallographic assignment requires an electronic-structure EFG calculation and comparison with the complete experimental spectrum.

An additional conditional scale test defines $P_s=|G_{ZZ,s}|\sqrt{1+\eta_{{\rm pc},s}^2/3}$ and fits one common positive multiplier to the four high-frequency NQR values. All 24 assignments to Cu1, Cu2, Cu4, and Cu5 are retained for each charge scenario; Cu3 is withheld from the calibration. The best frequency-only match in the uniform-charge model places the scaled Cu3 response at 1.503~MHz, close to the lower end of the archived low-frequency interval. The corresponding best matches across the six scenarios span 1.40--3.52~MHz.

This empirical scaling assumes a common conversion for all five sites and is not an independent prediction of an absolute nuclear EFG. Even the baseline match leaves a high-frequency root-mean-square residual of 0.584~MHz, with a maximum relative residual of 6.07\%, and does not reproduce the NMR-derived asymmetries. The comparison therefore motivates a targeted low-frequency assignment test without selecting an established charge distribution or a Cu-site mapping. All 144 assignments and their frequency and asymmetry discrepancies are retained in the supporting calculations.
